% CP-VI-0055
% target_chapter: VI/03 — Coordinate Rings
% solution_provenance: FRESH_CANONICAL_AUTHORED_SOLUTION

\subsection*{CP-VI-0055: Products and coordinate rings}

For affine algebraic sets $X\subseteq\mathbb A_k^m$ and $Y\subseteq\mathbb A_k^n$, show that
\[
k[X\times Y]\cong k[X]\otimes_k k[Y].
\]

\medskip
\noindent\textbf{Solution.}

Write $R=k[x_1,\dots,x_m]$, $S=k[y_1,\dots,y_n]$, and let
$J=I(X)\,k[x_1,\dots,x_m,y_1,\dots,y_n]
 +I(Y)\,k[x_1,\dots,x_m,y_1,\dots,y_n]$.
The standard tensor-product presentation gives
\[
k[X]\otimes_k k[Y]\cong
(R/I(X))\otimes_k(S/I(Y))\cong (R\otimes_k S)/J.
\]
We prove that $J=I(X\times Y)$, not merely that these ideals have
identical zero loci. The inclusion $J\subseteq I(X\times Y)$ is clear.
For the reverse inclusion let $H\in I(X\times Y)$. In
$R\otimes_k (S/I(Y))$ write its image as
\[
\overline H=\sum_{i=1}^r p_i(x)\otimes\overline q_i(y),
\]
where the $\overline q_i$ are linearly independent over $k$;
such an expression exists by choosing a basis for the finite-dimensional
span of the coefficient functions of $H$.
For each $a\in X$, the polynomial $H(a,y)$ vanishes at every point of
$Y$, hence its class in $S/I(Y)$ is zero. Therefore
\[
\sum_{i=1}^r p_i(a)\,\overline q_i=0.
\]
Linear independence gives $p_i(a)=0$ for all $i$ and all $a\in X$,
so $p_i\in I(X)$. Consequently $H\in J$.
Thus $J=I(X\times Y)$ and
\[
k[X\times Y]\cong k[X]\otimes_k k[Y].
\]
